% ---------------------------------------------------------------------------
% Preamble snippet — add to your main document (before \begin{document}):
%
%   \usepackage{listings}
%   \usepackage{xcolor}
%   \usepackage{booktabs}
%   \usepackage{hyperref}
%
%   \definecolor{codegreen}{rgb}{0,0.45,0.35}
%   \definecolor{codegray}{rgb}{0.45,0.45,0.45}
%
%   \lstdefinestyle{pythonstyle}{
%     language=Python,
%     basicstyle=\ttfamily\small,
%     keywordstyle=\bfseries,
%     commentstyle=\itshape\color{codegray},
%     stringstyle=\color{codegreen},
%     showstringspaces=false,
%     breaklines=true,
%     frame=single,
%     columns=fullflexible,
%     keepspaces=true,
%     literate={_}{{\_}}1,
%   }
%   \lstset{style=pythonstyle}
%
% Include this section with:
%   % ---------------------------------------------------------------------------
% Preamble snippet — add to your main document (before \begin{document}):
%
%   \usepackage{listings}
%   \usepackage{xcolor}
%   \usepackage{booktabs}
%   \usepackage{hyperref}
%
%   \definecolor{codegreen}{rgb}{0,0.45,0.35}
%   \definecolor{codegray}{rgb}{0.45,0.45,0.45}
%
%   \lstdefinestyle{pythonstyle}{
%     language=Python,
%     basicstyle=\ttfamily\small,
%     keywordstyle=\bfseries,
%     commentstyle=\itshape\color{codegray},
%     stringstyle=\color{codegreen},
%     showstringspaces=false,
%     breaklines=true,
%     frame=single,
%     columns=fullflexible,
%     keepspaces=true,
%     literate={_}{{\_}}1,
%   }
%   \lstset{style=pythonstyle}
%
% Include this section with:
%   % ---------------------------------------------------------------------------
% Preamble snippet — add to your main document (before \begin{document}):
%
%   \usepackage{listings}
%   \usepackage{xcolor}
%   \usepackage{booktabs}
%   \usepackage{hyperref}
%
%   \definecolor{codegreen}{rgb}{0,0.45,0.35}
%   \definecolor{codegray}{rgb}{0.45,0.45,0.45}
%
%   \lstdefinestyle{pythonstyle}{
%     language=Python,
%     basicstyle=\ttfamily\small,
%     keywordstyle=\bfseries,
%     commentstyle=\itshape\color{codegray},
%     stringstyle=\color{codegreen},
%     showstringspaces=false,
%     breaklines=true,
%     frame=single,
%     columns=fullflexible,
%     keepspaces=true,
%     literate={_}{{\_}}1,
%   }
%   \lstset{style=pythonstyle}
%
% Include this section with:
%   % ---------------------------------------------------------------------------
% Preamble snippet — add to your main document (before \begin{document}):
%
%   \usepackage{listings}
%   \usepackage{xcolor}
%   \usepackage{booktabs}
%   \usepackage{hyperref}
%
%   \definecolor{codegreen}{rgb}{0,0.45,0.35}
%   \definecolor{codegray}{rgb}{0.45,0.45,0.45}
%
%   \lstdefinestyle{pythonstyle}{
%     language=Python,
%     basicstyle=\ttfamily\small,
%     keywordstyle=\bfseries,
%     commentstyle=\itshape\color{codegray},
%     stringstyle=\color{codegreen},
%     showstringspaces=false,
%     breaklines=true,
%     frame=single,
%     columns=fullflexible,
%     keepspaces=true,
%     literate={_}{{\_}}1,
%   }
%   \lstset{style=pythonstyle}
%
% Include this section with:
%   \input{docs/n2p2_workflow_section}
% ---------------------------------------------------------------------------

\section{n2p2 workflow}
\label{sec:n2p2-workflow}

The \texttt{aiida-n2p2} plugin (version~0.4.0) provides an AiiDA-native interface to the
n2p2 neural-network potential package.
The workflow layer follows AiiDA best practices: each pipeline stage
is exposed as an independent \texttt{WorkChain}, and a composite work chain orchestrates
the full potential-building procedure.
Version~0.4.0 adds automatic restart when a training CalcJob stops before reaching the
target number of epochs (e.g.\ after scheduler walltime limits), while continuing to
merge learning curves across segments.
This section summarises the architecture, the workflow \emph{input API}, memory-conscious
retrieval defaults, restart behaviour, and the regression test suite.

\subsection{Workflow architecture}

The end-to-end pipeline consists of three physical steps:
\begin{enumerate}
  \item \textbf{Scaling} (\texttt{nnp-scaling}): generation of symmetry-function
        scaling data from the training set;
  \item \textbf{Training} (\texttt{nnp-train}): optimisation of neural-network
        weights using the scaled descriptors;
  \item \textbf{Validation} (optional, via \texttt{aiida-lammps}): a short LAMMPS
        molecular-dynamics run to verify the trained potential.
\end{enumerate}

Each step is wrapped in a dedicated work chain and can be submitted independently or
composed through \texttt{MakeNNPWorkchain}:

\begin{table}[htbp]
  \centering
  \caption{AiiDA workflow entry points registered by \texttt{aiida-n2p2} v0.4.0.}
  \label{tab:n2p2-workflows}
  \begin{tabular}{lll}
    \toprule
    Entry point & Class & Purpose \\
    \midrule
    \texttt{n2p2.scale}            & \texttt{N2p2ScaleWorkChain}            & Run scaling \\
    \texttt{n2p2.train}            & \texttt{N2p2TrainWorkChain}            & Run training (with restart) \\
    \texttt{n2p2.validate\_lammps} & \texttt{N2p2LammpsValidationWorkChain} & LAMMPS validation \\
    \texttt{n2p2.prepare\_inputs}  & \texttt{N2p2PrepareInputsWorkChain}    & Build \texttt{input.data}/\texttt{input.nn} \\
    \texttt{n2p2.make\_potential}  & \texttt{MakeNNPWorkchain}              & Full pipeline \\
    \bottomrule
  \end{tabular}
\end{table}

\texttt{MakeNNPWorkchain} exposes inputs from the three sub-work chains under the
namespaces \texttt{scaling}, \texttt{training}, and \texttt{validation}.
Scaling outputs (\texttt{inputData}, \texttt{inputNN}, \texttt{scale}) are forwarded
automatically to the training step; trained weights and network files are forwarded to
the validation step when enabled.
Workflow outputs are \texttt{potential} (best weights file), \texttt{scale}
(scaling data), \texttt{last\_weights} (checkpoint for restart),
\texttt{learning\_curve\_merged}, and \texttt{learning\_curve\_plot\_data}.

Validation is controlled by the top-level boolean input \texttt{run\_validation}
(default: \texttt{True}).
When set to \texttt{False}, the work chain stops after training.

% ---------------------------------------------------------------------------
\subsection{The workflow input API}
\label{subsec:input-api}

\subsubsection{What is an input API?}

In software engineering, an \textbf{API} (\emph{Application Programming Interface})
is the set of rules that defines \emph{how} a user or another program interacts with
a piece of software: which functions can be called, which arguments they accept, and
what they return.
For example, a web service might expose an API with endpoints such as
\texttt{GET /users/\{id\}}; clients must follow that contract to obtain data.

In AiiDA, a \textbf{workflow input API} plays an analogous role.
Each \texttt{WorkChain} declares a fixed list of \textbf{input ports}---named slots
that accept AiiDA data nodes (e.g.\ \texttt{Code}, \texttt{SinglefileData},
\texttt{Int}, \texttt{Dict}).
Together, these ports form the \emph{input schema} of the workflow: they specify
\begin{itemize}
  \item \textbf{which} quantities must be provided (e.g.\ training data, executable,
        atomic number);
  \item \textbf{under which names} they must be passed (e.g.\ \texttt{inputData}
        rather than \texttt{data\_file});
  \item \textbf{which AiiDA type} each input must have (e.g.\ \texttt{Int} for the
        atomic number);
  \item whether an input is \textbf{required or optional}, and its \textbf{default}
        value if optional;
  \item how inputs are \textbf{grouped} into \textbf{namespaces} when a workflow
        composes several sub-workflows (e.g.\ \texttt{scaling.code},
        \texttt{training.code}).
\end{itemize}

From the user's perspective, launching a workflow means building a Python dictionary
(\texttt{inputs}) whose keys and nested structure must match this schema exactly.
AiiDA validates the dictionary before submission; if a required port is missing or
has the wrong type, the submission fails immediately.
This validation is what makes the input API a stable \emph{contract} between the
plugin developer and the end user.

When a plugin is updated, the input API may change: ports can be renamed, regrouped,
or moved between namespaces.
Such changes are called \textbf{breaking changes} if existing scripts must be
rewritten to work with the new version.
Documenting the input API is therefore essential for reproducibility and for
migrating calculation scripts across plugin releases.

\subsubsection{Input API of \texttt{MakeNNPWorkchain} (v0.2.0)}

Version~0.2.0 refactors the input API of \texttt{MakeNNPWorkchain}.
The previous single nested namespace \texttt{n2p2} was replaced by three top-level
namespaces---\texttt{scaling}, \texttt{training}, and \texttt{validation}---plus the
boolean flag \texttt{run\_validation}.
Port names in the validation step were also aligned with AiiDA conventions
(\texttt{script}, \texttt{structure}).

\paragraph{Previous input API (pre-0.2.0).}
\begin{lstlisting}
inputs = {
    "n2p2": {
        "scale": {
            "code": scale_code,
            "nbin": nbin,
            "inputData": input_data,
            "inputNN": input_nn,
            "metadata": scale_metadata,
        },
        "train": {
            "code": train_code,
            "atomicNumber": atomic_number,
            "metadata": train_metadata,
        },
        "validate": {
            "code": lammps_code,
            "lammpsScript": lammps_script,
            "lammpsData": lammps_structure,
            "metadata": validate_metadata,
        },
    }
}
\end{lstlisting}

\paragraph{Current input API (v0.2.0).}
\begin{lstlisting}
from aiida.orm import Bool

inputs = {
    "scaling": {
        "code": scale_code,
        "nbin": nbin,
        "inputData": input_data,
        "inputNN": input_nn,
        "metadata": scale_metadata,
    },
    "training": {
        "code": train_code,
        "atomicNumber": atomic_number,  # e.g. Int(13) for Al
        "metadata": train_metadata,
    },
    "validation": {
        "code": lammps_code,
        "script": lammps_script,        # was: lammpsScript
        "structure": lammps_structure,  # was: lammpsData
        "metadata": validate_metadata,
        "retrieve_trajectory": Bool(False),
    },
    "run_validation": Bool(True),
}
\end{lstlisting}

Table~\ref{tab:input-api-changes} summarises the main differences between the two
versions.

\begin{table}[htbp]
  \centering
  \caption{Summary of input API changes in \texttt{aiida-n2p2} v0.2.0.}
  \label{tab:input-api-changes}
  \begin{tabular}{lll}
    \toprule
    Aspect & Pre-0.2.0 & v0.2.0 \\
    \midrule
    Top-level namespace   & \texttt{n2p2}     & \texttt{scaling}, \texttt{training}, \texttt{validation} \\
    Scaling sub-namespace & \texttt{scale}    & \texttt{scaling} \\
    Training sub-namespace& \texttt{train}    & \texttt{training} \\
    Validation sub-namespace & \texttt{validate} & \texttt{validation} \\
    LAMMPS input script   & \texttt{lammpsScript} & \texttt{script} \\
    LAMMPS structure file & \texttt{lammpsData}   & \texttt{structure} \\
    Skip validation flag  & not available     & \texttt{run\_validation} \\
    Trajectory retrieval  & always on         & \texttt{retrieve\_trajectory} (default \texttt{False}) \\
    \bottomrule
  \end{tabular}
\end{table}

The atomic number (\texttt{atomicNumber}) is a \emph{required} training input in
both versions.
It is not hard-coded in the plugin: the user must supply the correct value for each
chemical element (e.g.\ \texttt{Int(13)} for aluminium, \texttt{Int(5)} for boron).

To submit the workflow once the \texttt{inputs} dictionary is built:
\begin{lstlisting}
from aiida.engine import run_get_node
from aiida.plugins import WorkflowFactory

MakeNNP = WorkflowFactory("n2p2.make_potential")
result, node = run_get_node(MakeNNP, **inputs)
print(f"Submitted MakeNNPWorkchain with PK: {node.pk}")
\end{lstlisting}

% ---------------------------------------------------------------------------
\subsection{CalcJobs and parsers}

Three \texttt{CalcJob} plugins mirror the n2p2 executables:

\begin{table}[htbp]
  \centering
  \caption{Calculation and parser entry points.}
  \label{tab:n2p2-calcs}
  \begin{tabular}{lll}
    \toprule
    Entry point & CalcJob & Executable \\
    \midrule
    \texttt{n2p2.scale}   & \texttt{nnpScaling}  & \texttt{nnp-scaling} \\
    \texttt{n2p2.train}   & \texttt{nnpTraining} & \texttt{nnp-train}   \\
    \texttt{n2p2.predict} & \texttt{nnpPredict}  & \texttt{nnp-predict} \\
    \bottomrule
  \end{tabular}
\end{table}

CalcJobs were hardened for AiiDA compliance:
\begin{itemize}
  \item explicit \texttt{code} input in \texttt{define()};
  \item default scheduler resources and parser names;
  \item structured exit codes for missing output files;
  \item removal of debug \texttt{print()} statements from parsers;
  \item clearer error messages in scaling and predict parsers.
\end{itemize}

The training parser selects the best epoch from \texttt{learning-curve.out} (lowest
test-set RMSE) and stores the corresponding weight file as \texttt{weights}.
It also stores \texttt{last\_weights} from the final completed epoch, which serves
as the restart checkpoint for the next training session.

\paragraph{Remote output recovery (v0.4.0).}
On some HPC schedulers (notably OAR on Dahu), the default retrieve step can copy
incomplete output files before the executable finishes.
The scale and train parsers therefore poll the remote folder when needed:
\texttt{scaling.data} and \texttt{learning-curve.out} are read only after the file
size stabilises; training may prefer remote weight files when the learning curve
lags behind the last written checkpoint.
This allows automatic restart to proceed even when a CalcJob exits with a non-zero
scheduler status but parseable partial outputs exist.

% ---------------------------------------------------------------------------
\subsection{Structured inputs and restart (v0.3.0)}
\label{subsec:v030-restart}

\subsubsection{Data types}

Two AiiDA data types wrap n2p2 input files with provenance-friendly metadata:

\begin{table}[htbp]
  \centering
  \caption{Structured data entry points (v0.3.0).}
  \label{tab:n2p2-data}
  \begin{tabular}{lll}
    \toprule
    Entry point & Class & Purpose \\
    \midrule
    \texttt{n2p2.dataset}     & \texttt{N2p2Dataset}     & Training set (\texttt{input.data}) \\
    \texttt{n2p2.parameters}  & \texttt{N2p2Parameters}  & Network config (\texttt{input.nn}) \\
    \bottomrule
  \end{tabular}
\end{table}

\texttt{N2p2Parameters} parses \texttt{input.nn} into an editable keyword dictionary.
Keywords can be active, commented out, or flag-style (e.g.\
\texttt{\#use\_old\_weights\_short}).
The method \texttt{prepare\_for\_restart()} enables \texttt{use\_old\_weights\_short}
and optionally extends \texttt{epochs}:
\begin{lstlisting}
from aiida.plugins import DataFactory

N2p2Parameters = DataFactory("n2p2.parameters")
params = N2p2Parameters.from_file("input.nn")
params.set("epochs", "20")
params.prepare_for_restart(additional_epochs=20)
input_nn = params.get_singlefile()
\end{lstlisting}

\texttt{N2p2PrepareInputsWorkChain} renders \texttt{input.data} and \texttt{input.nn}
from these nodes (or from legacy \texttt{SinglefileData}), applying optional keyword
overrides and restart preparation when \texttt{is\_restart = Bool(True)}.

\subsubsection{Two restart modes compared}
\label{subsubsec:restart-modes}

Training can be continued in two ways.
Both use the same CalcJob restart engine (\texttt{use\_old\_weights\_short},
\texttt{weights.ZZZ.data}, merged learning curves), but differ in orchestration.

\begin{table}[htbp]
  \centering
  \caption{Manual vs.\ automatic training restart in \texttt{N2p2TrainWorkChain}.}
  \label{tab:restart-modes}
  \begin{tabular}{lll}
    \toprule
    Aspect & Manual (v0.3) & Automatic (v0.4) \\
    \midrule
    Inputs & \texttt{is\_restart}, \texttt{previous\_session} & \texttt{auto\_restart}, \texttt{max\_auto\_restarts} \\
    Who triggers next segment & User (new \texttt{submit}) & WorkChain (\texttt{while\_} loop) \\
    Typical use & Deliberate multi-step training & HPC walltime limits \\
    Prior segment status & Usually \texttt{finished\_ok} & May be failed (e.g.\ exit 143) \\
    Parseable outputs required & Yes & Yes (\texttt{learning-curve.out}, \texttt{last\_weights}) \\
    AiiDA provenance & N WorkChain nodes & 1 WorkChain, N CalcJob children \\
    Extend \texttt{epochs} & \texttt{additional\_epochs} (remaining in segment) & Fixed target in \texttt{input.nn} \\
    Example script & \texttt{restart\_train\_demo.py} & \texttt{auto\_restart\_train\_demo.py} \\
    \bottomrule
  \end{tabular}
\end{table}

\paragraph{Manual restart (inter-workflow).}
The user submits one WorkChain per session and links them with
\texttt{previous\_session}.
Each session is normally expected to finish successfully.
This is appropriate when you deliberately split training (e.g.\ 40~+~40 epochs
toward a target of~80) or when you want full control between steps.
Each manual session is a separate \texttt{N2p2TrainWorkChain} submission.

\paragraph{Automatic restart (intra-workflow).}
The user submits \emph{once}.
If a CalcJob stops before \texttt{epochs} is reached in \texttt{input.nn}
(e.g.\ scheduler walltime) but the parser recovered
\texttt{learning-curve.out} and checkpoint weights, the WorkChain submits
another CalcJob automatically until the target is met or
\texttt{max\_auto\_restarts} is exceeded.
This is appropriate on HPC clusters with fixed job time limits.

\paragraph{Output files and overwrite safety.}
Each CalcJob segment runs in a \emph{new} remote folder (new AiiDA UUID).
Segments do not overwrite each other's cluster files.
Intermediate segments remain as child calculations in the provenance graph.
The WorkChain output \texttt{learning\_curve\_merged} combines all segments;
\texttt{weights} and \texttt{last\_weights} refer to the final segment.

\subsubsection{Manual restart API}
\label{subsubsec:manual-restart}

\texttt{N2p2TrainWorkChain} accepts optional restart inputs:

\begin{itemize}
  \item \texttt{is\_restart} (\texttt{Bool}): restart from a previous session;
  \item \texttt{previous\_session} (\texttt{WorkChainNode}, \texttt{non\_db=True}):
        prior \texttt{N2p2TrainWorkChain} whose \texttt{last\_weights} are reused;
  \item \texttt{parameters} (\texttt{N2p2Parameters}): template used to regenerate
        \texttt{input.nn} with restart keywords;
  \item \texttt{additional\_epochs} (\texttt{Int}): number of \emph{remaining}
        epochs written into the restart segment's \texttt{input.nn} (not added to
        the cumulative total in the file).
\end{itemize}

On manual restart the WorkChain sets \texttt{epochs} in the regenerated
\texttt{input.nn} to \texttt{additional\_epochs} after enabling
\texttt{use\_old\_weights\_short}.
The global target for completion checks comes from \texttt{parameters} (e.g.\
\texttt{epochs=80} while the first segment used~40).
Automatic restart uses the same ``remaining epochs'' convention internally.

On restart, the work chain copies \texttt{previous\_session.outputs.last\_weights}
into the CalcJob sandbox as \texttt{weights.ZZZ.data}, sets \texttt{run\_label} to
the previous value plus one, and merges learning-curve segments into
\texttt{learning\_curve\_merged} and \texttt{learning\_curve\_plot\_data}
(one colour per run, global epoch axis).

Example (two deliberate sessions, 40~+~40 epochs toward target~80):
\begin{lstlisting}
from aiida.engine import run_get_node
from aiida.orm import Bool, Int
from aiida.plugins import WorkflowFactory

ScaleWC = WorkflowFactory("n2p2.scale")
TrainWC = WorkflowFactory("n2p2.train")

# Use examples/2.HPC/input.data and input.nn (Al dataset)
params.set("epochs", "40")
_, scaling = run_get_node(ScaleWC, ...)
_, run1 = run_get_node(TrainWC, ..., parameters=params)

params_restart = params.clone()
params_restart.set("epochs", "80")  # global target for WorkChain
_, run2 = run_get_node(
    TrainWC,
    ...,
    is_restart=Bool(True),
    previous_session=run1,
    parameters=params_restart,
    additional_epochs=Int(40),  # remaining epochs in segment 2
)
merged = run2.outputs.learning_curve_plot_data.get_dict()
assert merged["n_runs"] == 2
\end{lstlisting}

\subsubsection{Automatic restart API (v0.4.0)}
\label{subsubsec:auto-restart}

Additional inputs:

\begin{itemize}
  \item \texttt{auto\_restart} (\texttt{Bool}): loop until target \texttt{epochs};
  \item \texttt{max\_auto\_restarts} (\texttt{Int}): safety limit (default 10).
\end{itemize}

Set \texttt{max\_wallclock\_seconds} in CalcJob metadata so each segment is
stopped by the scheduler before the target is reached.
The train parser records \texttt{target\_epochs}, \texttt{training\_completed},
and \texttt{epochs\_remaining} in \texttt{training\_summary}.

\begin{lstlisting}
from aiida.engine import run_get_node
from aiida.orm import Bool, Int
from aiida.plugins import WorkflowFactory

TrainWC = WorkflowFactory("n2p2.train")

_, node = run_get_node(
    TrainWC,
    ...,
    parameters=params,
    auto_restart=Bool(True),
    max_auto_restarts=Int(10),
    metadata=Dict(dict={
        "options": {
            "resources": {"num_machines": 1},
            "max_wallclock_seconds": 7200,
        }
    }),
)
assert node.outputs.training_summary.get_dict()["training_completed"]
assert node.outputs.learning_curve_plot_data.get_dict()["n_runs"] >= 2
\end{lstlisting}

Walltime is enforced by the cluster scheduler; on a plain local AiiDA computer
the job may complete in one segment and automatic restart will not trigger.
Use \texttt{examples/3.Restart/auto\_restart\_dahu.py} on Dahu/OAR
(default \texttt{max\_wallclock\_seconds=360}, target~300 epochs) or
\texttt{auto\_restart\_train\_demo.py} on other schedulers to verify
end-to-end behaviour (\texttt{n\_runs >= 2} in \texttt{learning\_curve\_plot\_data}).
Logic is also covered by unit tests without re-running n2p2.

Template leftovers from the original plugin scaffold (\texttt{helpers.py},
\texttt{cli.py}, unused \texttt{DiffParameters} data type) were removed.

\subsection{Memory-conscious retrieval defaults}

Two design choices reduce permanent storage in the AiiDA repository and avoid
unnecessary data transfer on HPC systems.

\paragraph{Training weight files.}
n2p2 writes one \texttt{weights.ZZZ.EEEEEE.out} file per training epoch.
Only the file corresponding to the best epoch is needed as a workflow output.
Accordingly, in \texttt{nnpTraining}:
\begin{itemize}
  \item \texttt{learning-curve.out} is listed in \texttt{retrieve\_list}
        (stored permanently);
  \item \texttt{weights.*.out} is listed in \texttt{retrieve\_temporary\_list}
        (discarded after parsing).
\end{itemize}
The train parser reads the best weight file from the temporary retrieved folder and
stores only that single file in the provenance graph.
For a typical aluminium run with ${\sim}200$ epochs, this avoids archiving
${\sim}200$ redundant weight files (${\sim}1.2\,\mathrm{MB}$ in the reference case,
scaling to hundreds of megabytes for long training runs).

\paragraph{LAMMPS trajectories.}
By default, LAMMPS validation does \emph{not} retrieve trajectory files
(\texttt{*.lammpstrj}), which can grow to gigabytes for long molecular-dynamics runs.
Retrieval is opt-in via \texttt{validation.retrieve\_trajectory = Bool(True)}.
Examples that post-process trajectories (e.g.\ radial distribution function analysis
or structure visualisation) set this flag to \texttt{True}; production validation
runs that only check exit status can leave it at the default \texttt{False}.

\subsection{Regression tests}

A fast regression test suite (48 tests, ${\sim}6\,\mathrm{s}$) validates parser
logic, structured inputs, restart wiring, and learning-curve merging against golden
outputs from a successful aluminium run
(\texttt{MakeNNPWorkchain<949>}) without re-executing n2p2 or LAMMPS.
Fixtures are stored under \texttt{tests/fixtures/al/} with a
\texttt{REFERENCE.json} manifest recording:
\begin{itemize}
  \item MD5 checksum of \texttt{scaling.data};
  \item best training epoch (182), last epoch (200), and test-set RMSE;
  \item MD5 checksum of the selected best and last weights files.
\end{itemize}

Tests cover scale and train parser output, learning-curve epoch selection, plugin
entry-point registration, work-chain port specifications, \texttt{input.nn} parsing,
\texttt{N2p2Parameters} restart preparation, and CalcJob/WorkChain restart wiring.
They are run with:
\begin{lstlisting}[language=bash]
PYTHONPATH=. pytest -v
\end{lstlisting}

\subsection{Updated examples}

The following examples reflect the v0.4.0 API:
\begin{itemize}
  \item \texttt{examples/1.Al/AiidA-n2p2\_demo.ipynb} --- local Al demonstration
        notebook (with \texttt{retrieve\_trajectory = Bool(True)} for trajectory
        visualisation);
  \item \texttt{examples/2.HPC/wkchain\_Al.py} --- full pipeline on HPC or
        localhost (\texttt{N2P2\_COMPUTER}, code labels via environment variables);
  \item \texttt{examples/1.Boron/aiida-n2p2\_demo.py} --- boron example with
        \texttt{Int(5)} and \texttt{retrieve\_trajectory = Bool(False)};
  \item \texttt{examples/3.Restart/restart\_train\_demo.py} --- manual restart:
        two WorkChain submits (default 40~+~40 epochs; Al inputs from
        \texttt{examples/2.HPC/});
  \item \texttt{examples/3.Restart/auto\_restart\_dahu.py} --- automatic restart
        on Dahu/OAR (\texttt{pr-diamond}, configurable walltime per segment);
  \item \texttt{examples/3.Restart/auto\_restart\_train\_demo.py} --- automatic
        restart smoke test (any scheduler; localhost walltime may not trigger
        multiple segments).
\end{itemize}

Restart examples require the full Al \texttt{input.data} from
\texttt{examples/2.HPC/}; the minimal files under \texttt{examples/3.Restart/} are
for unit tests only.

% ---------------------------------------------------------------------------

\section{n2p2 workflow}
\label{sec:n2p2-workflow}

The \texttt{aiida-n2p2} plugin (version~0.4.0) provides an AiiDA-native interface to the
n2p2 neural-network potential package.
The workflow layer follows AiiDA best practices: each pipeline stage
is exposed as an independent \texttt{WorkChain}, and a composite work chain orchestrates
the full potential-building procedure.
Version~0.4.0 adds automatic restart when a training CalcJob stops before reaching the
target number of epochs (e.g.\ after scheduler walltime limits), while continuing to
merge learning curves across segments.
This section summarises the architecture, the workflow \emph{input API}, memory-conscious
retrieval defaults, restart behaviour, and the regression test suite.

\subsection{Workflow architecture}

The end-to-end pipeline consists of three physical steps:
\begin{enumerate}
  \item \textbf{Scaling} (\texttt{nnp-scaling}): generation of symmetry-function
        scaling data from the training set;
  \item \textbf{Training} (\texttt{nnp-train}): optimisation of neural-network
        weights using the scaled descriptors;
  \item \textbf{Validation} (optional, via \texttt{aiida-lammps}): a short LAMMPS
        molecular-dynamics run to verify the trained potential.
\end{enumerate}

Each step is wrapped in a dedicated work chain and can be submitted independently or
composed through \texttt{MakeNNPWorkchain}:

\begin{table}[htbp]
  \centering
  \caption{AiiDA workflow entry points registered by \texttt{aiida-n2p2} v0.4.0.}
  \label{tab:n2p2-workflows}
  \begin{tabular}{lll}
    \toprule
    Entry point & Class & Purpose \\
    \midrule
    \texttt{n2p2.scale}            & \texttt{N2p2ScaleWorkChain}            & Run scaling \\
    \texttt{n2p2.train}            & \texttt{N2p2TrainWorkChain}            & Run training (with restart) \\
    \texttt{n2p2.validate\_lammps} & \texttt{N2p2LammpsValidationWorkChain} & LAMMPS validation \\
    \texttt{n2p2.prepare\_inputs}  & \texttt{N2p2PrepareInputsWorkChain}    & Build \texttt{input.data}/\texttt{input.nn} \\
    \texttt{n2p2.make\_potential}  & \texttt{MakeNNPWorkchain}              & Full pipeline \\
    \bottomrule
  \end{tabular}
\end{table}

\texttt{MakeNNPWorkchain} exposes inputs from the three sub-work chains under the
namespaces \texttt{scaling}, \texttt{training}, and \texttt{validation}.
Scaling outputs (\texttt{inputData}, \texttt{inputNN}, \texttt{scale}) are forwarded
automatically to the training step; trained weights and network files are forwarded to
the validation step when enabled.
Workflow outputs are \texttt{potential} (best weights file), \texttt{scale}
(scaling data), \texttt{last\_weights} (checkpoint for restart),
\texttt{learning\_curve\_merged}, and \texttt{learning\_curve\_plot\_data}.

Validation is controlled by the top-level boolean input \texttt{run\_validation}
(default: \texttt{True}).
When set to \texttt{False}, the work chain stops after training.

% ---------------------------------------------------------------------------
\subsection{The workflow input API}
\label{subsec:input-api}

\subsubsection{What is an input API?}

In software engineering, an \textbf{API} (\emph{Application Programming Interface})
is the set of rules that defines \emph{how} a user or another program interacts with
a piece of software: which functions can be called, which arguments they accept, and
what they return.
For example, a web service might expose an API with endpoints such as
\texttt{GET /users/\{id\}}; clients must follow that contract to obtain data.

In AiiDA, a \textbf{workflow input API} plays an analogous role.
Each \texttt{WorkChain} declares a fixed list of \textbf{input ports}---named slots
that accept AiiDA data nodes (e.g.\ \texttt{Code}, \texttt{SinglefileData},
\texttt{Int}, \texttt{Dict}).
Together, these ports form the \emph{input schema} of the workflow: they specify
\begin{itemize}
  \item \textbf{which} quantities must be provided (e.g.\ training data, executable,
        atomic number);
  \item \textbf{under which names} they must be passed (e.g.\ \texttt{inputData}
        rather than \texttt{data\_file});
  \item \textbf{which AiiDA type} each input must have (e.g.\ \texttt{Int} for the
        atomic number);
  \item whether an input is \textbf{required or optional}, and its \textbf{default}
        value if optional;
  \item how inputs are \textbf{grouped} into \textbf{namespaces} when a workflow
        composes several sub-workflows (e.g.\ \texttt{scaling.code},
        \texttt{training.code}).
\end{itemize}

From the user's perspective, launching a workflow means building a Python dictionary
(\texttt{inputs}) whose keys and nested structure must match this schema exactly.
AiiDA validates the dictionary before submission; if a required port is missing or
has the wrong type, the submission fails immediately.
This validation is what makes the input API a stable \emph{contract} between the
plugin developer and the end user.

When a plugin is updated, the input API may change: ports can be renamed, regrouped,
or moved between namespaces.
Such changes are called \textbf{breaking changes} if existing scripts must be
rewritten to work with the new version.
Documenting the input API is therefore essential for reproducibility and for
migrating calculation scripts across plugin releases.

\subsubsection{Input API of \texttt{MakeNNPWorkchain} (v0.2.0)}

Version~0.2.0 refactors the input API of \texttt{MakeNNPWorkchain}.
The previous single nested namespace \texttt{n2p2} was replaced by three top-level
namespaces---\texttt{scaling}, \texttt{training}, and \texttt{validation}---plus the
boolean flag \texttt{run\_validation}.
Port names in the validation step were also aligned with AiiDA conventions
(\texttt{script}, \texttt{structure}).

\paragraph{Previous input API (pre-0.2.0).}
\begin{lstlisting}
inputs = {
    "n2p2": {
        "scale": {
            "code": scale_code,
            "nbin": nbin,
            "inputData": input_data,
            "inputNN": input_nn,
            "metadata": scale_metadata,
        },
        "train": {
            "code": train_code,
            "atomicNumber": atomic_number,
            "metadata": train_metadata,
        },
        "validate": {
            "code": lammps_code,
            "lammpsScript": lammps_script,
            "lammpsData": lammps_structure,
            "metadata": validate_metadata,
        },
    }
}
\end{lstlisting}

\paragraph{Current input API (v0.2.0).}
\begin{lstlisting}
from aiida.orm import Bool

inputs = {
    "scaling": {
        "code": scale_code,
        "nbin": nbin,
        "inputData": input_data,
        "inputNN": input_nn,
        "metadata": scale_metadata,
    },
    "training": {
        "code": train_code,
        "atomicNumber": atomic_number,  # e.g. Int(13) for Al
        "metadata": train_metadata,
    },
    "validation": {
        "code": lammps_code,
        "script": lammps_script,        # was: lammpsScript
        "structure": lammps_structure,  # was: lammpsData
        "metadata": validate_metadata,
        "retrieve_trajectory": Bool(False),
    },
    "run_validation": Bool(True),
}
\end{lstlisting}

Table~\ref{tab:input-api-changes} summarises the main differences between the two
versions.

\begin{table}[htbp]
  \centering
  \caption{Summary of input API changes in \texttt{aiida-n2p2} v0.2.0.}
  \label{tab:input-api-changes}
  \begin{tabular}{lll}
    \toprule
    Aspect & Pre-0.2.0 & v0.2.0 \\
    \midrule
    Top-level namespace   & \texttt{n2p2}     & \texttt{scaling}, \texttt{training}, \texttt{validation} \\
    Scaling sub-namespace & \texttt{scale}    & \texttt{scaling} \\
    Training sub-namespace& \texttt{train}    & \texttt{training} \\
    Validation sub-namespace & \texttt{validate} & \texttt{validation} \\
    LAMMPS input script   & \texttt{lammpsScript} & \texttt{script} \\
    LAMMPS structure file & \texttt{lammpsData}   & \texttt{structure} \\
    Skip validation flag  & not available     & \texttt{run\_validation} \\
    Trajectory retrieval  & always on         & \texttt{retrieve\_trajectory} (default \texttt{False}) \\
    \bottomrule
  \end{tabular}
\end{table}

The atomic number (\texttt{atomicNumber}) is a \emph{required} training input in
both versions.
It is not hard-coded in the plugin: the user must supply the correct value for each
chemical element (e.g.\ \texttt{Int(13)} for aluminium, \texttt{Int(5)} for boron).

To submit the workflow once the \texttt{inputs} dictionary is built:
\begin{lstlisting}
from aiida.engine import run_get_node
from aiida.plugins import WorkflowFactory

MakeNNP = WorkflowFactory("n2p2.make_potential")
result, node = run_get_node(MakeNNP, **inputs)
print(f"Submitted MakeNNPWorkchain with PK: {node.pk}")
\end{lstlisting}

% ---------------------------------------------------------------------------
\subsection{CalcJobs and parsers}

Three \texttt{CalcJob} plugins mirror the n2p2 executables:

\begin{table}[htbp]
  \centering
  \caption{Calculation and parser entry points.}
  \label{tab:n2p2-calcs}
  \begin{tabular}{lll}
    \toprule
    Entry point & CalcJob & Executable \\
    \midrule
    \texttt{n2p2.scale}   & \texttt{nnpScaling}  & \texttt{nnp-scaling} \\
    \texttt{n2p2.train}   & \texttt{nnpTraining} & \texttt{nnp-train}   \\
    \texttt{n2p2.predict} & \texttt{nnpPredict}  & \texttt{nnp-predict} \\
    \bottomrule
  \end{tabular}
\end{table}

CalcJobs were hardened for AiiDA compliance:
\begin{itemize}
  \item explicit \texttt{code} input in \texttt{define()};
  \item default scheduler resources and parser names;
  \item structured exit codes for missing output files;
  \item removal of debug \texttt{print()} statements from parsers;
  \item clearer error messages in scaling and predict parsers.
\end{itemize}

The training parser selects the best epoch from \texttt{learning-curve.out} (lowest
test-set RMSE) and stores the corresponding weight file as \texttt{weights}.
It also stores \texttt{last\_weights} from the final completed epoch, which serves
as the restart checkpoint for the next training session.

\paragraph{Remote output recovery (v0.4.0).}
On some HPC schedulers (notably OAR on Dahu), the default retrieve step can copy
incomplete output files before the executable finishes.
The scale and train parsers therefore poll the remote folder when needed:
\texttt{scaling.data} and \texttt{learning-curve.out} are read only after the file
size stabilises; training may prefer remote weight files when the learning curve
lags behind the last written checkpoint.
This allows automatic restart to proceed even when a CalcJob exits with a non-zero
scheduler status but parseable partial outputs exist.

% ---------------------------------------------------------------------------
\subsection{Structured inputs and restart (v0.3.0)}
\label{subsec:v030-restart}

\subsubsection{Data types}

Two AiiDA data types wrap n2p2 input files with provenance-friendly metadata:

\begin{table}[htbp]
  \centering
  \caption{Structured data entry points (v0.3.0).}
  \label{tab:n2p2-data}
  \begin{tabular}{lll}
    \toprule
    Entry point & Class & Purpose \\
    \midrule
    \texttt{n2p2.dataset}     & \texttt{N2p2Dataset}     & Training set (\texttt{input.data}) \\
    \texttt{n2p2.parameters}  & \texttt{N2p2Parameters}  & Network config (\texttt{input.nn}) \\
    \bottomrule
  \end{tabular}
\end{table}

\texttt{N2p2Parameters} parses \texttt{input.nn} into an editable keyword dictionary.
Keywords can be active, commented out, or flag-style (e.g.\
\texttt{\#use\_old\_weights\_short}).
The method \texttt{prepare\_for\_restart()} enables \texttt{use\_old\_weights\_short}
and optionally extends \texttt{epochs}:
\begin{lstlisting}
from aiida.plugins import DataFactory

N2p2Parameters = DataFactory("n2p2.parameters")
params = N2p2Parameters.from_file("input.nn")
params.set("epochs", "20")
params.prepare_for_restart(additional_epochs=20)
input_nn = params.get_singlefile()
\end{lstlisting}

\texttt{N2p2PrepareInputsWorkChain} renders \texttt{input.data} and \texttt{input.nn}
from these nodes (or from legacy \texttt{SinglefileData}), applying optional keyword
overrides and restart preparation when \texttt{is\_restart = Bool(True)}.

\subsubsection{Two restart modes compared}
\label{subsubsec:restart-modes}

Training can be continued in two ways.
Both use the same CalcJob restart engine (\texttt{use\_old\_weights\_short},
\texttt{weights.ZZZ.data}, merged learning curves), but differ in orchestration.

\begin{table}[htbp]
  \centering
  \caption{Manual vs.\ automatic training restart in \texttt{N2p2TrainWorkChain}.}
  \label{tab:restart-modes}
  \begin{tabular}{lll}
    \toprule
    Aspect & Manual (v0.3) & Automatic (v0.4) \\
    \midrule
    Inputs & \texttt{is\_restart}, \texttt{previous\_session} & \texttt{auto\_restart}, \texttt{max\_auto\_restarts} \\
    Who triggers next segment & User (new \texttt{submit}) & WorkChain (\texttt{while\_} loop) \\
    Typical use & Deliberate multi-step training & HPC walltime limits \\
    Prior segment status & Usually \texttt{finished\_ok} & May be failed (e.g.\ exit 143) \\
    Parseable outputs required & Yes & Yes (\texttt{learning-curve.out}, \texttt{last\_weights}) \\
    AiiDA provenance & N WorkChain nodes & 1 WorkChain, N CalcJob children \\
    Extend \texttt{epochs} & \texttt{additional\_epochs} (remaining in segment) & Fixed target in \texttt{input.nn} \\
    Example script & \texttt{restart\_train\_demo.py} & \texttt{auto\_restart\_train\_demo.py} \\
    \bottomrule
  \end{tabular}
\end{table}

\paragraph{Manual restart (inter-workflow).}
The user submits one WorkChain per session and links them with
\texttt{previous\_session}.
Each session is normally expected to finish successfully.
This is appropriate when you deliberately split training (e.g.\ 40~+~40 epochs
toward a target of~80) or when you want full control between steps.
Each manual session is a separate \texttt{N2p2TrainWorkChain} submission.

\paragraph{Automatic restart (intra-workflow).}
The user submits \emph{once}.
If a CalcJob stops before \texttt{epochs} is reached in \texttt{input.nn}
(e.g.\ scheduler walltime) but the parser recovered
\texttt{learning-curve.out} and checkpoint weights, the WorkChain submits
another CalcJob automatically until the target is met or
\texttt{max\_auto\_restarts} is exceeded.
This is appropriate on HPC clusters with fixed job time limits.

\paragraph{Output files and overwrite safety.}
Each CalcJob segment runs in a \emph{new} remote folder (new AiiDA UUID).
Segments do not overwrite each other's cluster files.
Intermediate segments remain as child calculations in the provenance graph.
The WorkChain output \texttt{learning\_curve\_merged} combines all segments;
\texttt{weights} and \texttt{last\_weights} refer to the final segment.

\subsubsection{Manual restart API}
\label{subsubsec:manual-restart}

\texttt{N2p2TrainWorkChain} accepts optional restart inputs:

\begin{itemize}
  \item \texttt{is\_restart} (\texttt{Bool}): restart from a previous session;
  \item \texttt{previous\_session} (\texttt{WorkChainNode}, \texttt{non\_db=True}):
        prior \texttt{N2p2TrainWorkChain} whose \texttt{last\_weights} are reused;
  \item \texttt{parameters} (\texttt{N2p2Parameters}): template used to regenerate
        \texttt{input.nn} with restart keywords;
  \item \texttt{additional\_epochs} (\texttt{Int}): number of \emph{remaining}
        epochs written into the restart segment's \texttt{input.nn} (not added to
        the cumulative total in the file).
\end{itemize}

On manual restart the WorkChain sets \texttt{epochs} in the regenerated
\texttt{input.nn} to \texttt{additional\_epochs} after enabling
\texttt{use\_old\_weights\_short}.
The global target for completion checks comes from \texttt{parameters} (e.g.\
\texttt{epochs=80} while the first segment used~40).
Automatic restart uses the same ``remaining epochs'' convention internally.

On restart, the work chain copies \texttt{previous\_session.outputs.last\_weights}
into the CalcJob sandbox as \texttt{weights.ZZZ.data}, sets \texttt{run\_label} to
the previous value plus one, and merges learning-curve segments into
\texttt{learning\_curve\_merged} and \texttt{learning\_curve\_plot\_data}
(one colour per run, global epoch axis).

Example (two deliberate sessions, 40~+~40 epochs toward target~80):
\begin{lstlisting}
from aiida.engine import run_get_node
from aiida.orm import Bool, Int
from aiida.plugins import WorkflowFactory

ScaleWC = WorkflowFactory("n2p2.scale")
TrainWC = WorkflowFactory("n2p2.train")

# Use examples/2.HPC/input.data and input.nn (Al dataset)
params.set("epochs", "40")
_, scaling = run_get_node(ScaleWC, ...)
_, run1 = run_get_node(TrainWC, ..., parameters=params)

params_restart = params.clone()
params_restart.set("epochs", "80")  # global target for WorkChain
_, run2 = run_get_node(
    TrainWC,
    ...,
    is_restart=Bool(True),
    previous_session=run1,
    parameters=params_restart,
    additional_epochs=Int(40),  # remaining epochs in segment 2
)
merged = run2.outputs.learning_curve_plot_data.get_dict()
assert merged["n_runs"] == 2
\end{lstlisting}

\subsubsection{Automatic restart API (v0.4.0)}
\label{subsubsec:auto-restart}

Additional inputs:

\begin{itemize}
  \item \texttt{auto\_restart} (\texttt{Bool}): loop until target \texttt{epochs};
  \item \texttt{max\_auto\_restarts} (\texttt{Int}): safety limit (default 10).
\end{itemize}

Set \texttt{max\_wallclock\_seconds} in CalcJob metadata so each segment is
stopped by the scheduler before the target is reached.
The train parser records \texttt{target\_epochs}, \texttt{training\_completed},
and \texttt{epochs\_remaining} in \texttt{training\_summary}.

\begin{lstlisting}
from aiida.engine import run_get_node
from aiida.orm import Bool, Int
from aiida.plugins import WorkflowFactory

TrainWC = WorkflowFactory("n2p2.train")

_, node = run_get_node(
    TrainWC,
    ...,
    parameters=params,
    auto_restart=Bool(True),
    max_auto_restarts=Int(10),
    metadata=Dict(dict={
        "options": {
            "resources": {"num_machines": 1},
            "max_wallclock_seconds": 7200,
        }
    }),
)
assert node.outputs.training_summary.get_dict()["training_completed"]
assert node.outputs.learning_curve_plot_data.get_dict()["n_runs"] >= 2
\end{lstlisting}

Walltime is enforced by the cluster scheduler; on a plain local AiiDA computer
the job may complete in one segment and automatic restart will not trigger.
Use \texttt{examples/3.Restart/auto\_restart\_dahu.py} on Dahu/OAR
(default \texttt{max\_wallclock\_seconds=360}, target~300 epochs) or
\texttt{auto\_restart\_train\_demo.py} on other schedulers to verify
end-to-end behaviour (\texttt{n\_runs >= 2} in \texttt{learning\_curve\_plot\_data}).
Logic is also covered by unit tests without re-running n2p2.

Template leftovers from the original plugin scaffold (\texttt{helpers.py},
\texttt{cli.py}, unused \texttt{DiffParameters} data type) were removed.

\subsection{Memory-conscious retrieval defaults}

Two design choices reduce permanent storage in the AiiDA repository and avoid
unnecessary data transfer on HPC systems.

\paragraph{Training weight files.}
n2p2 writes one \texttt{weights.ZZZ.EEEEEE.out} file per training epoch.
Only the file corresponding to the best epoch is needed as a workflow output.
Accordingly, in \texttt{nnpTraining}:
\begin{itemize}
  \item \texttt{learning-curve.out} is listed in \texttt{retrieve\_list}
        (stored permanently);
  \item \texttt{weights.*.out} is listed in \texttt{retrieve\_temporary\_list}
        (discarded after parsing).
\end{itemize}
The train parser reads the best weight file from the temporary retrieved folder and
stores only that single file in the provenance graph.
For a typical aluminium run with ${\sim}200$ epochs, this avoids archiving
${\sim}200$ redundant weight files (${\sim}1.2\,\mathrm{MB}$ in the reference case,
scaling to hundreds of megabytes for long training runs).

\paragraph{LAMMPS trajectories.}
By default, LAMMPS validation does \emph{not} retrieve trajectory files
(\texttt{*.lammpstrj}), which can grow to gigabytes for long molecular-dynamics runs.
Retrieval is opt-in via \texttt{validation.retrieve\_trajectory = Bool(True)}.
Examples that post-process trajectories (e.g.\ radial distribution function analysis
or structure visualisation) set this flag to \texttt{True}; production validation
runs that only check exit status can leave it at the default \texttt{False}.

\subsection{Regression tests}

A fast regression test suite (48 tests, ${\sim}6\,\mathrm{s}$) validates parser
logic, structured inputs, restart wiring, and learning-curve merging against golden
outputs from a successful aluminium run
(\texttt{MakeNNPWorkchain<949>}) without re-executing n2p2 or LAMMPS.
Fixtures are stored under \texttt{tests/fixtures/al/} with a
\texttt{REFERENCE.json} manifest recording:
\begin{itemize}
  \item MD5 checksum of \texttt{scaling.data};
  \item best training epoch (182), last epoch (200), and test-set RMSE;
  \item MD5 checksum of the selected best and last weights files.
\end{itemize}

Tests cover scale and train parser output, learning-curve epoch selection, plugin
entry-point registration, work-chain port specifications, \texttt{input.nn} parsing,
\texttt{N2p2Parameters} restart preparation, and CalcJob/WorkChain restart wiring.
They are run with:
\begin{lstlisting}[language=bash]
PYTHONPATH=. pytest -v
\end{lstlisting}

\subsection{Updated examples}

The following examples reflect the v0.4.0 API:
\begin{itemize}
  \item \texttt{examples/1.Al/AiidA-n2p2\_demo.ipynb} --- local Al demonstration
        notebook (with \texttt{retrieve\_trajectory = Bool(True)} for trajectory
        visualisation);
  \item \texttt{examples/2.HPC/wkchain\_Al.py} --- full pipeline on HPC or
        localhost (\texttt{N2P2\_COMPUTER}, code labels via environment variables);
  \item \texttt{examples/1.Boron/aiida-n2p2\_demo.py} --- boron example with
        \texttt{Int(5)} and \texttt{retrieve\_trajectory = Bool(False)};
  \item \texttt{examples/3.Restart/restart\_train\_demo.py} --- manual restart:
        two WorkChain submits (default 40~+~40 epochs; Al inputs from
        \texttt{examples/2.HPC/});
  \item \texttt{examples/3.Restart/auto\_restart\_dahu.py} --- automatic restart
        on Dahu/OAR (\texttt{pr-diamond}, configurable walltime per segment);
  \item \texttt{examples/3.Restart/auto\_restart\_train\_demo.py} --- automatic
        restart smoke test (any scheduler; localhost walltime may not trigger
        multiple segments).
\end{itemize}

Restart examples require the full Al \texttt{input.data} from
\texttt{examples/2.HPC/}; the minimal files under \texttt{examples/3.Restart/} are
for unit tests only.

% ---------------------------------------------------------------------------

\section{n2p2 workflow}
\label{sec:n2p2-workflow}

The \texttt{aiida-n2p2} plugin (version~0.4.0) provides an AiiDA-native interface to the
n2p2 neural-network potential package.
The workflow layer follows AiiDA best practices: each pipeline stage
is exposed as an independent \texttt{WorkChain}, and a composite work chain orchestrates
the full potential-building procedure.
Version~0.4.0 adds automatic restart when a training CalcJob stops before reaching the
target number of epochs (e.g.\ after scheduler walltime limits), while continuing to
merge learning curves across segments.
This section summarises the architecture, the workflow \emph{input API}, memory-conscious
retrieval defaults, restart behaviour, and the regression test suite.

\subsection{Workflow architecture}

The end-to-end pipeline consists of three physical steps:
\begin{enumerate}
  \item \textbf{Scaling} (\texttt{nnp-scaling}): generation of symmetry-function
        scaling data from the training set;
  \item \textbf{Training} (\texttt{nnp-train}): optimisation of neural-network
        weights using the scaled descriptors;
  \item \textbf{Validation} (optional, via \texttt{aiida-lammps}): a short LAMMPS
        molecular-dynamics run to verify the trained potential.
\end{enumerate}

Each step is wrapped in a dedicated work chain and can be submitted independently or
composed through \texttt{MakeNNPWorkchain}:

\begin{table}[htbp]
  \centering
  \caption{AiiDA workflow entry points registered by \texttt{aiida-n2p2} v0.4.0.}
  \label{tab:n2p2-workflows}
  \begin{tabular}{lll}
    \toprule
    Entry point & Class & Purpose \\
    \midrule
    \texttt{n2p2.scale}            & \texttt{N2p2ScaleWorkChain}            & Run scaling \\
    \texttt{n2p2.train}            & \texttt{N2p2TrainWorkChain}            & Run training (with restart) \\
    \texttt{n2p2.validate\_lammps} & \texttt{N2p2LammpsValidationWorkChain} & LAMMPS validation \\
    \texttt{n2p2.prepare\_inputs}  & \texttt{N2p2PrepareInputsWorkChain}    & Build \texttt{input.data}/\texttt{input.nn} \\
    \texttt{n2p2.make\_potential}  & \texttt{MakeNNPWorkchain}              & Full pipeline \\
    \bottomrule
  \end{tabular}
\end{table}

\texttt{MakeNNPWorkchain} exposes inputs from the three sub-work chains under the
namespaces \texttt{scaling}, \texttt{training}, and \texttt{validation}.
Scaling outputs (\texttt{inputData}, \texttt{inputNN}, \texttt{scale}) are forwarded
automatically to the training step; trained weights and network files are forwarded to
the validation step when enabled.
Workflow outputs are \texttt{potential} (best weights file), \texttt{scale}
(scaling data), \texttt{last\_weights} (checkpoint for restart),
\texttt{learning\_curve\_merged}, and \texttt{learning\_curve\_plot\_data}.

Validation is controlled by the top-level boolean input \texttt{run\_validation}
(default: \texttt{True}).
When set to \texttt{False}, the work chain stops after training.

% ---------------------------------------------------------------------------
\subsection{The workflow input API}
\label{subsec:input-api}

\subsubsection{What is an input API?}

In software engineering, an \textbf{API} (\emph{Application Programming Interface})
is the set of rules that defines \emph{how} a user or another program interacts with
a piece of software: which functions can be called, which arguments they accept, and
what they return.
For example, a web service might expose an API with endpoints such as
\texttt{GET /users/\{id\}}; clients must follow that contract to obtain data.

In AiiDA, a \textbf{workflow input API} plays an analogous role.
Each \texttt{WorkChain} declares a fixed list of \textbf{input ports}---named slots
that accept AiiDA data nodes (e.g.\ \texttt{Code}, \texttt{SinglefileData},
\texttt{Int}, \texttt{Dict}).
Together, these ports form the \emph{input schema} of the workflow: they specify
\begin{itemize}
  \item \textbf{which} quantities must be provided (e.g.\ training data, executable,
        atomic number);
  \item \textbf{under which names} they must be passed (e.g.\ \texttt{inputData}
        rather than \texttt{data\_file});
  \item \textbf{which AiiDA type} each input must have (e.g.\ \texttt{Int} for the
        atomic number);
  \item whether an input is \textbf{required or optional}, and its \textbf{default}
        value if optional;
  \item how inputs are \textbf{grouped} into \textbf{namespaces} when a workflow
        composes several sub-workflows (e.g.\ \texttt{scaling.code},
        \texttt{training.code}).
\end{itemize}

From the user's perspective, launching a workflow means building a Python dictionary
(\texttt{inputs}) whose keys and nested structure must match this schema exactly.
AiiDA validates the dictionary before submission; if a required port is missing or
has the wrong type, the submission fails immediately.
This validation is what makes the input API a stable \emph{contract} between the
plugin developer and the end user.

When a plugin is updated, the input API may change: ports can be renamed, regrouped,
or moved between namespaces.
Such changes are called \textbf{breaking changes} if existing scripts must be
rewritten to work with the new version.
Documenting the input API is therefore essential for reproducibility and for
migrating calculation scripts across plugin releases.

\subsubsection{Input API of \texttt{MakeNNPWorkchain} (v0.2.0)}

Version~0.2.0 refactors the input API of \texttt{MakeNNPWorkchain}.
The previous single nested namespace \texttt{n2p2} was replaced by three top-level
namespaces---\texttt{scaling}, \texttt{training}, and \texttt{validation}---plus the
boolean flag \texttt{run\_validation}.
Port names in the validation step were also aligned with AiiDA conventions
(\texttt{script}, \texttt{structure}).

\paragraph{Previous input API (pre-0.2.0).}
\begin{lstlisting}
inputs = {
    "n2p2": {
        "scale": {
            "code": scale_code,
            "nbin": nbin,
            "inputData": input_data,
            "inputNN": input_nn,
            "metadata": scale_metadata,
        },
        "train": {
            "code": train_code,
            "atomicNumber": atomic_number,
            "metadata": train_metadata,
        },
        "validate": {
            "code": lammps_code,
            "lammpsScript": lammps_script,
            "lammpsData": lammps_structure,
            "metadata": validate_metadata,
        },
    }
}
\end{lstlisting}

\paragraph{Current input API (v0.2.0).}
\begin{lstlisting}
from aiida.orm import Bool

inputs = {
    "scaling": {
        "code": scale_code,
        "nbin": nbin,
        "inputData": input_data,
        "inputNN": input_nn,
        "metadata": scale_metadata,
    },
    "training": {
        "code": train_code,
        "atomicNumber": atomic_number,  # e.g. Int(13) for Al
        "metadata": train_metadata,
    },
    "validation": {
        "code": lammps_code,
        "script": lammps_script,        # was: lammpsScript
        "structure": lammps_structure,  # was: lammpsData
        "metadata": validate_metadata,
        "retrieve_trajectory": Bool(False),
    },
    "run_validation": Bool(True),
}
\end{lstlisting}

Table~\ref{tab:input-api-changes} summarises the main differences between the two
versions.

\begin{table}[htbp]
  \centering
  \caption{Summary of input API changes in \texttt{aiida-n2p2} v0.2.0.}
  \label{tab:input-api-changes}
  \begin{tabular}{lll}
    \toprule
    Aspect & Pre-0.2.0 & v0.2.0 \\
    \midrule
    Top-level namespace   & \texttt{n2p2}     & \texttt{scaling}, \texttt{training}, \texttt{validation} \\
    Scaling sub-namespace & \texttt{scale}    & \texttt{scaling} \\
    Training sub-namespace& \texttt{train}    & \texttt{training} \\
    Validation sub-namespace & \texttt{validate} & \texttt{validation} \\
    LAMMPS input script   & \texttt{lammpsScript} & \texttt{script} \\
    LAMMPS structure file & \texttt{lammpsData}   & \texttt{structure} \\
    Skip validation flag  & not available     & \texttt{run\_validation} \\
    Trajectory retrieval  & always on         & \texttt{retrieve\_trajectory} (default \texttt{False}) \\
    \bottomrule
  \end{tabular}
\end{table}

The atomic number (\texttt{atomicNumber}) is a \emph{required} training input in
both versions.
It is not hard-coded in the plugin: the user must supply the correct value for each
chemical element (e.g.\ \texttt{Int(13)} for aluminium, \texttt{Int(5)} for boron).

To submit the workflow once the \texttt{inputs} dictionary is built:
\begin{lstlisting}
from aiida.engine import run_get_node
from aiida.plugins import WorkflowFactory

MakeNNP = WorkflowFactory("n2p2.make_potential")
result, node = run_get_node(MakeNNP, **inputs)
print(f"Submitted MakeNNPWorkchain with PK: {node.pk}")
\end{lstlisting}

% ---------------------------------------------------------------------------
\subsection{CalcJobs and parsers}

Three \texttt{CalcJob} plugins mirror the n2p2 executables:

\begin{table}[htbp]
  \centering
  \caption{Calculation and parser entry points.}
  \label{tab:n2p2-calcs}
  \begin{tabular}{lll}
    \toprule
    Entry point & CalcJob & Executable \\
    \midrule
    \texttt{n2p2.scale}   & \texttt{nnpScaling}  & \texttt{nnp-scaling} \\
    \texttt{n2p2.train}   & \texttt{nnpTraining} & \texttt{nnp-train}   \\
    \texttt{n2p2.predict} & \texttt{nnpPredict}  & \texttt{nnp-predict} \\
    \bottomrule
  \end{tabular}
\end{table}

CalcJobs were hardened for AiiDA compliance:
\begin{itemize}
  \item explicit \texttt{code} input in \texttt{define()};
  \item default scheduler resources and parser names;
  \item structured exit codes for missing output files;
  \item removal of debug \texttt{print()} statements from parsers;
  \item clearer error messages in scaling and predict parsers.
\end{itemize}

The training parser selects the best epoch from \texttt{learning-curve.out} (lowest
test-set RMSE) and stores the corresponding weight file as \texttt{weights}.
It also stores \texttt{last\_weights} from the final completed epoch, which serves
as the restart checkpoint for the next training session.

\paragraph{Remote output recovery (v0.4.0).}
On some HPC schedulers (notably OAR on Dahu), the default retrieve step can copy
incomplete output files before the executable finishes.
The scale and train parsers therefore poll the remote folder when needed:
\texttt{scaling.data} and \texttt{learning-curve.out} are read only after the file
size stabilises; training may prefer remote weight files when the learning curve
lags behind the last written checkpoint.
This allows automatic restart to proceed even when a CalcJob exits with a non-zero
scheduler status but parseable partial outputs exist.

% ---------------------------------------------------------------------------
\subsection{Structured inputs and restart (v0.3.0)}
\label{subsec:v030-restart}

\subsubsection{Data types}

Two AiiDA data types wrap n2p2 input files with provenance-friendly metadata:

\begin{table}[htbp]
  \centering
  \caption{Structured data entry points (v0.3.0).}
  \label{tab:n2p2-data}
  \begin{tabular}{lll}
    \toprule
    Entry point & Class & Purpose \\
    \midrule
    \texttt{n2p2.dataset}     & \texttt{N2p2Dataset}     & Training set (\texttt{input.data}) \\
    \texttt{n2p2.parameters}  & \texttt{N2p2Parameters}  & Network config (\texttt{input.nn}) \\
    \bottomrule
  \end{tabular}
\end{table}

\texttt{N2p2Parameters} parses \texttt{input.nn} into an editable keyword dictionary.
Keywords can be active, commented out, or flag-style (e.g.\
\texttt{\#use\_old\_weights\_short}).
The method \texttt{prepare\_for\_restart()} enables \texttt{use\_old\_weights\_short}
and optionally extends \texttt{epochs}:
\begin{lstlisting}
from aiida.plugins import DataFactory

N2p2Parameters = DataFactory("n2p2.parameters")
params = N2p2Parameters.from_file("input.nn")
params.set("epochs", "20")
params.prepare_for_restart(additional_epochs=20)
input_nn = params.get_singlefile()
\end{lstlisting}

\texttt{N2p2PrepareInputsWorkChain} renders \texttt{input.data} and \texttt{input.nn}
from these nodes (or from legacy \texttt{SinglefileData}), applying optional keyword
overrides and restart preparation when \texttt{is\_restart = Bool(True)}.

\subsubsection{Two restart modes compared}
\label{subsubsec:restart-modes}

Training can be continued in two ways.
Both use the same CalcJob restart engine (\texttt{use\_old\_weights\_short},
\texttt{weights.ZZZ.data}, merged learning curves), but differ in orchestration.

\begin{table}[htbp]
  \centering
  \caption{Manual vs.\ automatic training restart in \texttt{N2p2TrainWorkChain}.}
  \label{tab:restart-modes}
  \begin{tabular}{lll}
    \toprule
    Aspect & Manual (v0.3) & Automatic (v0.4) \\
    \midrule
    Inputs & \texttt{is\_restart}, \texttt{previous\_session} & \texttt{auto\_restart}, \texttt{max\_auto\_restarts} \\
    Who triggers next segment & User (new \texttt{submit}) & WorkChain (\texttt{while\_} loop) \\
    Typical use & Deliberate multi-step training & HPC walltime limits \\
    Prior segment status & Usually \texttt{finished\_ok} & May be failed (e.g.\ exit 143) \\
    Parseable outputs required & Yes & Yes (\texttt{learning-curve.out}, \texttt{last\_weights}) \\
    AiiDA provenance & N WorkChain nodes & 1 WorkChain, N CalcJob children \\
    Extend \texttt{epochs} & \texttt{additional\_epochs} (remaining in segment) & Fixed target in \texttt{input.nn} \\
    Example script & \texttt{restart\_train\_demo.py} & \texttt{auto\_restart\_train\_demo.py} \\
    \bottomrule
  \end{tabular}
\end{table}

\paragraph{Manual restart (inter-workflow).}
The user submits one WorkChain per session and links them with
\texttt{previous\_session}.
Each session is normally expected to finish successfully.
This is appropriate when you deliberately split training (e.g.\ 40~+~40 epochs
toward a target of~80) or when you want full control between steps.
Each manual session is a separate \texttt{N2p2TrainWorkChain} submission.

\paragraph{Automatic restart (intra-workflow).}
The user submits \emph{once}.
If a CalcJob stops before \texttt{epochs} is reached in \texttt{input.nn}
(e.g.\ scheduler walltime) but the parser recovered
\texttt{learning-curve.out} and checkpoint weights, the WorkChain submits
another CalcJob automatically until the target is met or
\texttt{max\_auto\_restarts} is exceeded.
This is appropriate on HPC clusters with fixed job time limits.

\paragraph{Output files and overwrite safety.}
Each CalcJob segment runs in a \emph{new} remote folder (new AiiDA UUID).
Segments do not overwrite each other's cluster files.
Intermediate segments remain as child calculations in the provenance graph.
The WorkChain output \texttt{learning\_curve\_merged} combines all segments;
\texttt{weights} and \texttt{last\_weights} refer to the final segment.

\subsubsection{Manual restart API}
\label{subsubsec:manual-restart}

\texttt{N2p2TrainWorkChain} accepts optional restart inputs:

\begin{itemize}
  \item \texttt{is\_restart} (\texttt{Bool}): restart from a previous session;
  \item \texttt{previous\_session} (\texttt{WorkChainNode}, \texttt{non\_db=True}):
        prior \texttt{N2p2TrainWorkChain} whose \texttt{last\_weights} are reused;
  \item \texttt{parameters} (\texttt{N2p2Parameters}): template used to regenerate
        \texttt{input.nn} with restart keywords;
  \item \texttt{additional\_epochs} (\texttt{Int}): number of \emph{remaining}
        epochs written into the restart segment's \texttt{input.nn} (not added to
        the cumulative total in the file).
\end{itemize}

On manual restart the WorkChain sets \texttt{epochs} in the regenerated
\texttt{input.nn} to \texttt{additional\_epochs} after enabling
\texttt{use\_old\_weights\_short}.
The global target for completion checks comes from \texttt{parameters} (e.g.\
\texttt{epochs=80} while the first segment used~40).
Automatic restart uses the same ``remaining epochs'' convention internally.

On restart, the work chain copies \texttt{previous\_session.outputs.last\_weights}
into the CalcJob sandbox as \texttt{weights.ZZZ.data}, sets \texttt{run\_label} to
the previous value plus one, and merges learning-curve segments into
\texttt{learning\_curve\_merged} and \texttt{learning\_curve\_plot\_data}
(one colour per run, global epoch axis).

Example (two deliberate sessions, 40~+~40 epochs toward target~80):
\begin{lstlisting}
from aiida.engine import run_get_node
from aiida.orm import Bool, Int
from aiida.plugins import WorkflowFactory

ScaleWC = WorkflowFactory("n2p2.scale")
TrainWC = WorkflowFactory("n2p2.train")

# Use examples/2.HPC/input.data and input.nn (Al dataset)
params.set("epochs", "40")
_, scaling = run_get_node(ScaleWC, ...)
_, run1 = run_get_node(TrainWC, ..., parameters=params)

params_restart = params.clone()
params_restart.set("epochs", "80")  # global target for WorkChain
_, run2 = run_get_node(
    TrainWC,
    ...,
    is_restart=Bool(True),
    previous_session=run1,
    parameters=params_restart,
    additional_epochs=Int(40),  # remaining epochs in segment 2
)
merged = run2.outputs.learning_curve_plot_data.get_dict()
assert merged["n_runs"] == 2
\end{lstlisting}

\subsubsection{Automatic restart API (v0.4.0)}
\label{subsubsec:auto-restart}

Additional inputs:

\begin{itemize}
  \item \texttt{auto\_restart} (\texttt{Bool}): loop until target \texttt{epochs};
  \item \texttt{max\_auto\_restarts} (\texttt{Int}): safety limit (default 10).
\end{itemize}

Set \texttt{max\_wallclock\_seconds} in CalcJob metadata so each segment is
stopped by the scheduler before the target is reached.
The train parser records \texttt{target\_epochs}, \texttt{training\_completed},
and \texttt{epochs\_remaining} in \texttt{training\_summary}.

\begin{lstlisting}
from aiida.engine import run_get_node
from aiida.orm import Bool, Int
from aiida.plugins import WorkflowFactory

TrainWC = WorkflowFactory("n2p2.train")

_, node = run_get_node(
    TrainWC,
    ...,
    parameters=params,
    auto_restart=Bool(True),
    max_auto_restarts=Int(10),
    metadata=Dict(dict={
        "options": {
            "resources": {"num_machines": 1},
            "max_wallclock_seconds": 7200,
        }
    }),
)
assert node.outputs.training_summary.get_dict()["training_completed"]
assert node.outputs.learning_curve_plot_data.get_dict()["n_runs"] >= 2
\end{lstlisting}

Walltime is enforced by the cluster scheduler; on a plain local AiiDA computer
the job may complete in one segment and automatic restart will not trigger.
Use \texttt{examples/3.Restart/auto\_restart\_dahu.py} on Dahu/OAR
(default \texttt{max\_wallclock\_seconds=360}, target~300 epochs) or
\texttt{auto\_restart\_train\_demo.py} on other schedulers to verify
end-to-end behaviour (\texttt{n\_runs >= 2} in \texttt{learning\_curve\_plot\_data}).
Logic is also covered by unit tests without re-running n2p2.

Template leftovers from the original plugin scaffold (\texttt{helpers.py},
\texttt{cli.py}, unused \texttt{DiffParameters} data type) were removed.

\subsection{Memory-conscious retrieval defaults}

Two design choices reduce permanent storage in the AiiDA repository and avoid
unnecessary data transfer on HPC systems.

\paragraph{Training weight files.}
n2p2 writes one \texttt{weights.ZZZ.EEEEEE.out} file per training epoch.
Only the file corresponding to the best epoch is needed as a workflow output.
Accordingly, in \texttt{nnpTraining}:
\begin{itemize}
  \item \texttt{learning-curve.out} is listed in \texttt{retrieve\_list}
        (stored permanently);
  \item \texttt{weights.*.out} is listed in \texttt{retrieve\_temporary\_list}
        (discarded after parsing).
\end{itemize}
The train parser reads the best weight file from the temporary retrieved folder and
stores only that single file in the provenance graph.
For a typical aluminium run with ${\sim}200$ epochs, this avoids archiving
${\sim}200$ redundant weight files (${\sim}1.2\,\mathrm{MB}$ in the reference case,
scaling to hundreds of megabytes for long training runs).

\paragraph{LAMMPS trajectories.}
By default, LAMMPS validation does \emph{not} retrieve trajectory files
(\texttt{*.lammpstrj}), which can grow to gigabytes for long molecular-dynamics runs.
Retrieval is opt-in via \texttt{validation.retrieve\_trajectory = Bool(True)}.
Examples that post-process trajectories (e.g.\ radial distribution function analysis
or structure visualisation) set this flag to \texttt{True}; production validation
runs that only check exit status can leave it at the default \texttt{False}.

\subsection{Regression tests}

A fast regression test suite (48 tests, ${\sim}6\,\mathrm{s}$) validates parser
logic, structured inputs, restart wiring, and learning-curve merging against golden
outputs from a successful aluminium run
(\texttt{MakeNNPWorkchain<949>}) without re-executing n2p2 or LAMMPS.
Fixtures are stored under \texttt{tests/fixtures/al/} with a
\texttt{REFERENCE.json} manifest recording:
\begin{itemize}
  \item MD5 checksum of \texttt{scaling.data};
  \item best training epoch (182), last epoch (200), and test-set RMSE;
  \item MD5 checksum of the selected best and last weights files.
\end{itemize}

Tests cover scale and train parser output, learning-curve epoch selection, plugin
entry-point registration, work-chain port specifications, \texttt{input.nn} parsing,
\texttt{N2p2Parameters} restart preparation, and CalcJob/WorkChain restart wiring.
They are run with:
\begin{lstlisting}[language=bash]
PYTHONPATH=. pytest -v
\end{lstlisting}

\subsection{Updated examples}

The following examples reflect the v0.4.0 API:
\begin{itemize}
  \item \texttt{examples/1.Al/AiidA-n2p2\_demo.ipynb} --- local Al demonstration
        notebook (with \texttt{retrieve\_trajectory = Bool(True)} for trajectory
        visualisation);
  \item \texttt{examples/2.HPC/wkchain\_Al.py} --- full pipeline on HPC or
        localhost (\texttt{N2P2\_COMPUTER}, code labels via environment variables);
  \item \texttt{examples/1.Boron/aiida-n2p2\_demo.py} --- boron example with
        \texttt{Int(5)} and \texttt{retrieve\_trajectory = Bool(False)};
  \item \texttt{examples/3.Restart/restart\_train\_demo.py} --- manual restart:
        two WorkChain submits (default 40~+~40 epochs; Al inputs from
        \texttt{examples/2.HPC/});
  \item \texttt{examples/3.Restart/auto\_restart\_dahu.py} --- automatic restart
        on Dahu/OAR (\texttt{pr-diamond}, configurable walltime per segment);
  \item \texttt{examples/3.Restart/auto\_restart\_train\_demo.py} --- automatic
        restart smoke test (any scheduler; localhost walltime may not trigger
        multiple segments).
\end{itemize}

Restart examples require the full Al \texttt{input.data} from
\texttt{examples/2.HPC/}; the minimal files under \texttt{examples/3.Restart/} are
for unit tests only.

% ---------------------------------------------------------------------------

\section{n2p2 workflow}
\label{sec:n2p2-workflow}

The \texttt{aiida-n2p2} plugin (version~0.4.0) provides an AiiDA-native interface to the
n2p2 neural-network potential package.
The workflow layer follows AiiDA best practices: each pipeline stage
is exposed as an independent \texttt{WorkChain}, and a composite work chain orchestrates
the full potential-building procedure.
Version~0.4.0 adds automatic restart when a training CalcJob stops before reaching the
target number of epochs (e.g.\ after scheduler walltime limits), while continuing to
merge learning curves across segments.
This section summarises the architecture, the workflow \emph{input API}, memory-conscious
retrieval defaults, restart behaviour, and the regression test suite.

\subsection{Workflow architecture}

The end-to-end pipeline consists of three physical steps:
\begin{enumerate}
  \item \textbf{Scaling} (\texttt{nnp-scaling}): generation of symmetry-function
        scaling data from the training set;
  \item \textbf{Training} (\texttt{nnp-train}): optimisation of neural-network
        weights using the scaled descriptors;
  \item \textbf{Validation} (optional, via \texttt{aiida-lammps}): a short LAMMPS
        molecular-dynamics run to verify the trained potential.
\end{enumerate}

Each step is wrapped in a dedicated work chain and can be submitted independently or
composed through \texttt{MakeNNPWorkchain}:

\begin{table}[htbp]
  \centering
  \caption{AiiDA workflow entry points registered by \texttt{aiida-n2p2} v0.4.0.}
  \label{tab:n2p2-workflows}
  \begin{tabular}{lll}
    \toprule
    Entry point & Class & Purpose \\
    \midrule
    \texttt{n2p2.scale}            & \texttt{N2p2ScaleWorkChain}            & Run scaling \\
    \texttt{n2p2.train}            & \texttt{N2p2TrainWorkChain}            & Run training (with restart) \\
    \texttt{n2p2.validate\_lammps} & \texttt{N2p2LammpsValidationWorkChain} & LAMMPS validation \\
    \texttt{n2p2.prepare\_inputs}  & \texttt{N2p2PrepareInputsWorkChain}    & Build \texttt{input.data}/\texttt{input.nn} \\
    \texttt{n2p2.make\_potential}  & \texttt{MakeNNPWorkchain}              & Full pipeline \\
    \bottomrule
  \end{tabular}
\end{table}

\texttt{MakeNNPWorkchain} exposes inputs from the three sub-work chains under the
namespaces \texttt{scaling}, \texttt{training}, and \texttt{validation}.
Scaling outputs (\texttt{inputData}, \texttt{inputNN}, \texttt{scale}) are forwarded
automatically to the training step; trained weights and network files are forwarded to
the validation step when enabled.
Workflow outputs are \texttt{potential} (best weights file), \texttt{scale}
(scaling data), \texttt{last\_weights} (checkpoint for restart),
\texttt{learning\_curve\_merged}, and \texttt{learning\_curve\_plot\_data}.

Validation is controlled by the top-level boolean input \texttt{run\_validation}
(default: \texttt{True}).
When set to \texttt{False}, the work chain stops after training.

% ---------------------------------------------------------------------------
\subsection{The workflow input API}
\label{subsec:input-api}

\subsubsection{What is an input API?}

In software engineering, an \textbf{API} (\emph{Application Programming Interface})
is the set of rules that defines \emph{how} a user or another program interacts with
a piece of software: which functions can be called, which arguments they accept, and
what they return.
For example, a web service might expose an API with endpoints such as
\texttt{GET /users/\{id\}}; clients must follow that contract to obtain data.

In AiiDA, a \textbf{workflow input API} plays an analogous role.
Each \texttt{WorkChain} declares a fixed list of \textbf{input ports}---named slots
that accept AiiDA data nodes (e.g.\ \texttt{Code}, \texttt{SinglefileData},
\texttt{Int}, \texttt{Dict}).
Together, these ports form the \emph{input schema} of the workflow: they specify
\begin{itemize}
  \item \textbf{which} quantities must be provided (e.g.\ training data, executable,
        atomic number);
  \item \textbf{under which names} they must be passed (e.g.\ \texttt{inputData}
        rather than \texttt{data\_file});
  \item \textbf{which AiiDA type} each input must have (e.g.\ \texttt{Int} for the
        atomic number);
  \item whether an input is \textbf{required or optional}, and its \textbf{default}
        value if optional;
  \item how inputs are \textbf{grouped} into \textbf{namespaces} when a workflow
        composes several sub-workflows (e.g.\ \texttt{scaling.code},
        \texttt{training.code}).
\end{itemize}

From the user's perspective, launching a workflow means building a Python dictionary
(\texttt{inputs}) whose keys and nested structure must match this schema exactly.
AiiDA validates the dictionary before submission; if a required port is missing or
has the wrong type, the submission fails immediately.
This validation is what makes the input API a stable \emph{contract} between the
plugin developer and the end user.

When a plugin is updated, the input API may change: ports can be renamed, regrouped,
or moved between namespaces.
Such changes are called \textbf{breaking changes} if existing scripts must be
rewritten to work with the new version.
Documenting the input API is therefore essential for reproducibility and for
migrating calculation scripts across plugin releases.

\subsubsection{Input API of \texttt{MakeNNPWorkchain} (v0.2.0)}

Version~0.2.0 refactors the input API of \texttt{MakeNNPWorkchain}.
The previous single nested namespace \texttt{n2p2} was replaced by three top-level
namespaces---\texttt{scaling}, \texttt{training}, and \texttt{validation}---plus the
boolean flag \texttt{run\_validation}.
Port names in the validation step were also aligned with AiiDA conventions
(\texttt{script}, \texttt{structure}).

\paragraph{Previous input API (pre-0.2.0).}
\begin{lstlisting}
inputs = {
    "n2p2": {
        "scale": {
            "code": scale_code,
            "nbin": nbin,
            "inputData": input_data,
            "inputNN": input_nn,
            "metadata": scale_metadata,
        },
        "train": {
            "code": train_code,
            "atomicNumber": atomic_number,
            "metadata": train_metadata,
        },
        "validate": {
            "code": lammps_code,
            "lammpsScript": lammps_script,
            "lammpsData": lammps_structure,
            "metadata": validate_metadata,
        },
    }
}
\end{lstlisting}

\paragraph{Current input API (v0.2.0).}
\begin{lstlisting}
from aiida.orm import Bool

inputs = {
    "scaling": {
        "code": scale_code,
        "nbin": nbin,
        "inputData": input_data,
        "inputNN": input_nn,
        "metadata": scale_metadata,
    },
    "training": {
        "code": train_code,
        "atomicNumber": atomic_number,  # e.g. Int(13) for Al
        "metadata": train_metadata,
    },
    "validation": {
        "code": lammps_code,
        "script": lammps_script,        # was: lammpsScript
        "structure": lammps_structure,  # was: lammpsData
        "metadata": validate_metadata,
        "retrieve_trajectory": Bool(False),
    },
    "run_validation": Bool(True),
}
\end{lstlisting}

Table~\ref{tab:input-api-changes} summarises the main differences between the two
versions.

\begin{table}[htbp]
  \centering
  \caption{Summary of input API changes in \texttt{aiida-n2p2} v0.2.0.}
  \label{tab:input-api-changes}
  \begin{tabular}{lll}
    \toprule
    Aspect & Pre-0.2.0 & v0.2.0 \\
    \midrule
    Top-level namespace   & \texttt{n2p2}     & \texttt{scaling}, \texttt{training}, \texttt{validation} \\
    Scaling sub-namespace & \texttt{scale}    & \texttt{scaling} \\
    Training sub-namespace& \texttt{train}    & \texttt{training} \\
    Validation sub-namespace & \texttt{validate} & \texttt{validation} \\
    LAMMPS input script   & \texttt{lammpsScript} & \texttt{script} \\
    LAMMPS structure file & \texttt{lammpsData}   & \texttt{structure} \\
    Skip validation flag  & not available     & \texttt{run\_validation} \\
    Trajectory retrieval  & always on         & \texttt{retrieve\_trajectory} (default \texttt{False}) \\
    \bottomrule
  \end{tabular}
\end{table}

The atomic number (\texttt{atomicNumber}) is a \emph{required} training input in
both versions.
It is not hard-coded in the plugin: the user must supply the correct value for each
chemical element (e.g.\ \texttt{Int(13)} for aluminium, \texttt{Int(5)} for boron).

To submit the workflow once the \texttt{inputs} dictionary is built:
\begin{lstlisting}
from aiida.engine import run_get_node
from aiida.plugins import WorkflowFactory

MakeNNP = WorkflowFactory("n2p2.make_potential")
result, node = run_get_node(MakeNNP, **inputs)
print(f"Submitted MakeNNPWorkchain with PK: {node.pk}")
\end{lstlisting}

% ---------------------------------------------------------------------------
\subsection{CalcJobs and parsers}

Three \texttt{CalcJob} plugins mirror the n2p2 executables:

\begin{table}[htbp]
  \centering
  \caption{Calculation and parser entry points.}
  \label{tab:n2p2-calcs}
  \begin{tabular}{lll}
    \toprule
    Entry point & CalcJob & Executable \\
    \midrule
    \texttt{n2p2.scale}   & \texttt{nnpScaling}  & \texttt{nnp-scaling} \\
    \texttt{n2p2.train}   & \texttt{nnpTraining} & \texttt{nnp-train}   \\
    \texttt{n2p2.predict} & \texttt{nnpPredict}  & \texttt{nnp-predict} \\
    \bottomrule
  \end{tabular}
\end{table}

CalcJobs were hardened for AiiDA compliance:
\begin{itemize}
  \item explicit \texttt{code} input in \texttt{define()};
  \item default scheduler resources and parser names;
  \item structured exit codes for missing output files;
  \item removal of debug \texttt{print()} statements from parsers;
  \item clearer error messages in scaling and predict parsers.
\end{itemize}

The training parser selects the best epoch from \texttt{learning-curve.out} (lowest
test-set RMSE) and stores the corresponding weight file as \texttt{weights}.
It also stores \texttt{last\_weights} from the final completed epoch, which serves
as the restart checkpoint for the next training session.

\paragraph{Remote output recovery (v0.4.0).}
On some HPC schedulers (notably OAR on Dahu), the default retrieve step can copy
incomplete output files before the executable finishes.
The scale and train parsers therefore poll the remote folder when needed:
\texttt{scaling.data} and \texttt{learning-curve.out} are read only after the file
size stabilises; training may prefer remote weight files when the learning curve
lags behind the last written checkpoint.
This allows automatic restart to proceed even when a CalcJob exits with a non-zero
scheduler status but parseable partial outputs exist.

% ---------------------------------------------------------------------------
\subsection{Structured inputs and restart (v0.3.0)}
\label{subsec:v030-restart}

\subsubsection{Data types}

Two AiiDA data types wrap n2p2 input files with provenance-friendly metadata:

\begin{table}[htbp]
  \centering
  \caption{Structured data entry points (v0.3.0).}
  \label{tab:n2p2-data}
  \begin{tabular}{lll}
    \toprule
    Entry point & Class & Purpose \\
    \midrule
    \texttt{n2p2.dataset}     & \texttt{N2p2Dataset}     & Training set (\texttt{input.data}) \\
    \texttt{n2p2.parameters}  & \texttt{N2p2Parameters}  & Network config (\texttt{input.nn}) \\
    \bottomrule
  \end{tabular}
\end{table}

\texttt{N2p2Parameters} parses \texttt{input.nn} into an editable keyword dictionary.
Keywords can be active, commented out, or flag-style (e.g.\
\texttt{\#use\_old\_weights\_short}).
The method \texttt{prepare\_for\_restart()} enables \texttt{use\_old\_weights\_short}
and optionally extends \texttt{epochs}:
\begin{lstlisting}
from aiida.plugins import DataFactory

N2p2Parameters = DataFactory("n2p2.parameters")
params = N2p2Parameters.from_file("input.nn")
params.set("epochs", "20")
params.prepare_for_restart(additional_epochs=20)
input_nn = params.get_singlefile()
\end{lstlisting}

\texttt{N2p2PrepareInputsWorkChain} renders \texttt{input.data} and \texttt{input.nn}
from these nodes (or from legacy \texttt{SinglefileData}), applying optional keyword
overrides and restart preparation when \texttt{is\_restart = Bool(True)}.

\subsubsection{Two restart modes compared}
\label{subsubsec:restart-modes}

Training can be continued in two ways.
Both use the same CalcJob restart engine (\texttt{use\_old\_weights\_short},
\texttt{weights.ZZZ.data}, merged learning curves), but differ in orchestration.

\begin{table}[htbp]
  \centering
  \caption{Manual vs.\ automatic training restart in \texttt{N2p2TrainWorkChain}.}
  \label{tab:restart-modes}
  \begin{tabular}{lll}
    \toprule
    Aspect & Manual (v0.3) & Automatic (v0.4) \\
    \midrule
    Inputs & \texttt{is\_restart}, \texttt{previous\_session} & \texttt{auto\_restart}, \texttt{max\_auto\_restarts} \\
    Who triggers next segment & User (new \texttt{submit}) & WorkChain (\texttt{while\_} loop) \\
    Typical use & Deliberate multi-step training & HPC walltime limits \\
    Prior segment status & Usually \texttt{finished\_ok} & May be failed (e.g.\ exit 143) \\
    Parseable outputs required & Yes & Yes (\texttt{learning-curve.out}, \texttt{last\_weights}) \\
    AiiDA provenance & N WorkChain nodes & 1 WorkChain, N CalcJob children \\
    Extend \texttt{epochs} & \texttt{additional\_epochs} (remaining in segment) & Fixed target in \texttt{input.nn} \\
    Example script & \texttt{restart\_train\_demo.py} & \texttt{auto\_restart\_train\_demo.py} \\
    \bottomrule
  \end{tabular}
\end{table}

\paragraph{Manual restart (inter-workflow).}
The user submits one WorkChain per session and links them with
\texttt{previous\_session}.
Each session is normally expected to finish successfully.
This is appropriate when you deliberately split training (e.g.\ 40~+~40 epochs
toward a target of~80) or when you want full control between steps.
Each manual session is a separate \texttt{N2p2TrainWorkChain} submission.

\paragraph{Automatic restart (intra-workflow).}
The user submits \emph{once}.
If a CalcJob stops before \texttt{epochs} is reached in \texttt{input.nn}
(e.g.\ scheduler walltime) but the parser recovered
\texttt{learning-curve.out} and checkpoint weights, the WorkChain submits
another CalcJob automatically until the target is met or
\texttt{max\_auto\_restarts} is exceeded.
This is appropriate on HPC clusters with fixed job time limits.

\paragraph{Output files and overwrite safety.}
Each CalcJob segment runs in a \emph{new} remote folder (new AiiDA UUID).
Segments do not overwrite each other's cluster files.
Intermediate segments remain as child calculations in the provenance graph.
The WorkChain output \texttt{learning\_curve\_merged} combines all segments;
\texttt{weights} and \texttt{last\_weights} refer to the final segment.

\subsubsection{Manual restart API}
\label{subsubsec:manual-restart}

\texttt{N2p2TrainWorkChain} accepts optional restart inputs:

\begin{itemize}
  \item \texttt{is\_restart} (\texttt{Bool}): restart from a previous session;
  \item \texttt{previous\_session} (\texttt{WorkChainNode}, \texttt{non\_db=True}):
        prior \texttt{N2p2TrainWorkChain} whose \texttt{last\_weights} are reused;
  \item \texttt{parameters} (\texttt{N2p2Parameters}): template used to regenerate
        \texttt{input.nn} with restart keywords;
  \item \texttt{additional\_epochs} (\texttt{Int}): number of \emph{remaining}
        epochs written into the restart segment's \texttt{input.nn} (not added to
        the cumulative total in the file).
\end{itemize}

On manual restart the WorkChain sets \texttt{epochs} in the regenerated
\texttt{input.nn} to \texttt{additional\_epochs} after enabling
\texttt{use\_old\_weights\_short}.
The global target for completion checks comes from \texttt{parameters} (e.g.\
\texttt{epochs=80} while the first segment used~40).
Automatic restart uses the same ``remaining epochs'' convention internally.

On restart, the work chain copies \texttt{previous\_session.outputs.last\_weights}
into the CalcJob sandbox as \texttt{weights.ZZZ.data}, sets \texttt{run\_label} to
the previous value plus one, and merges learning-curve segments into
\texttt{learning\_curve\_merged} and \texttt{learning\_curve\_plot\_data}
(one colour per run, global epoch axis).

Example (two deliberate sessions, 40~+~40 epochs toward target~80):
\begin{lstlisting}
from aiida.engine import run_get_node
from aiida.orm import Bool, Int
from aiida.plugins import WorkflowFactory

ScaleWC = WorkflowFactory("n2p2.scale")
TrainWC = WorkflowFactory("n2p2.train")

# Use examples/2.HPC/input.data and input.nn (Al dataset)
params.set("epochs", "40")
_, scaling = run_get_node(ScaleWC, ...)
_, run1 = run_get_node(TrainWC, ..., parameters=params)

params_restart = params.clone()
params_restart.set("epochs", "80")  # global target for WorkChain
_, run2 = run_get_node(
    TrainWC,
    ...,
    is_restart=Bool(True),
    previous_session=run1,
    parameters=params_restart,
    additional_epochs=Int(40),  # remaining epochs in segment 2
)
merged = run2.outputs.learning_curve_plot_data.get_dict()
assert merged["n_runs"] == 2
\end{lstlisting}

\subsubsection{Automatic restart API (v0.4.0)}
\label{subsubsec:auto-restart}

Additional inputs:

\begin{itemize}
  \item \texttt{auto\_restart} (\texttt{Bool}): loop until target \texttt{epochs};
  \item \texttt{max\_auto\_restarts} (\texttt{Int}): safety limit (default 10).
\end{itemize}

Set \texttt{max\_wallclock\_seconds} in CalcJob metadata so each segment is
stopped by the scheduler before the target is reached.
The train parser records \texttt{target\_epochs}, \texttt{training\_completed},
and \texttt{epochs\_remaining} in \texttt{training\_summary}.

\begin{lstlisting}
from aiida.engine import run_get_node
from aiida.orm import Bool, Int
from aiida.plugins import WorkflowFactory

TrainWC = WorkflowFactory("n2p2.train")

_, node = run_get_node(
    TrainWC,
    ...,
    parameters=params,
    auto_restart=Bool(True),
    max_auto_restarts=Int(10),
    metadata=Dict(dict={
        "options": {
            "resources": {"num_machines": 1},
            "max_wallclock_seconds": 7200,
        }
    }),
)
assert node.outputs.training_summary.get_dict()["training_completed"]
assert node.outputs.learning_curve_plot_data.get_dict()["n_runs"] >= 2
\end{lstlisting}

Walltime is enforced by the cluster scheduler; on a plain local AiiDA computer
the job may complete in one segment and automatic restart will not trigger.
Use \texttt{examples/3.Restart/auto\_restart\_dahu.py} on Dahu/OAR
(default \texttt{max\_wallclock\_seconds=360}, target~300 epochs) or
\texttt{auto\_restart\_train\_demo.py} on other schedulers to verify
end-to-end behaviour (\texttt{n\_runs >= 2} in \texttt{learning\_curve\_plot\_data}).
Logic is also covered by unit tests without re-running n2p2.

Template leftovers from the original plugin scaffold (\texttt{helpers.py},
\texttt{cli.py}, unused \texttt{DiffParameters} data type) were removed.

\subsection{Memory-conscious retrieval defaults}

Two design choices reduce permanent storage in the AiiDA repository and avoid
unnecessary data transfer on HPC systems.

\paragraph{Training weight files.}
n2p2 writes one \texttt{weights.ZZZ.EEEEEE.out} file per training epoch.
Only the file corresponding to the best epoch is needed as a workflow output.
Accordingly, in \texttt{nnpTraining}:
\begin{itemize}
  \item \texttt{learning-curve.out} is listed in \texttt{retrieve\_list}
        (stored permanently);
  \item \texttt{weights.*.out} is listed in \texttt{retrieve\_temporary\_list}
        (discarded after parsing).
\end{itemize}
The train parser reads the best weight file from the temporary retrieved folder and
stores only that single file in the provenance graph.
For a typical aluminium run with ${\sim}200$ epochs, this avoids archiving
${\sim}200$ redundant weight files (${\sim}1.2\,\mathrm{MB}$ in the reference case,
scaling to hundreds of megabytes for long training runs).

\paragraph{LAMMPS trajectories.}
By default, LAMMPS validation does \emph{not} retrieve trajectory files
(\texttt{*.lammpstrj}), which can grow to gigabytes for long molecular-dynamics runs.
Retrieval is opt-in via \texttt{validation.retrieve\_trajectory = Bool(True)}.
Examples that post-process trajectories (e.g.\ radial distribution function analysis
or structure visualisation) set this flag to \texttt{True}; production validation
runs that only check exit status can leave it at the default \texttt{False}.

\subsection{Regression tests}

A fast regression test suite (48 tests, ${\sim}6\,\mathrm{s}$) validates parser
logic, structured inputs, restart wiring, and learning-curve merging against golden
outputs from a successful aluminium run
(\texttt{MakeNNPWorkchain<949>}) without re-executing n2p2 or LAMMPS.
Fixtures are stored under \texttt{tests/fixtures/al/} with a
\texttt{REFERENCE.json} manifest recording:
\begin{itemize}
  \item MD5 checksum of \texttt{scaling.data};
  \item best training epoch (182), last epoch (200), and test-set RMSE;
  \item MD5 checksum of the selected best and last weights files.
\end{itemize}

Tests cover scale and train parser output, learning-curve epoch selection, plugin
entry-point registration, work-chain port specifications, \texttt{input.nn} parsing,
\texttt{N2p2Parameters} restart preparation, and CalcJob/WorkChain restart wiring.
They are run with:
\begin{lstlisting}[language=bash]
PYTHONPATH=. pytest -v
\end{lstlisting}

\subsection{Updated examples}

The following examples reflect the v0.4.0 API:
\begin{itemize}
  \item \texttt{examples/1.Al/AiidA-n2p2\_demo.ipynb} --- local Al demonstration
        notebook (with \texttt{retrieve\_trajectory = Bool(True)} for trajectory
        visualisation);
  \item \texttt{examples/2.HPC/wkchain\_Al.py} --- full pipeline on HPC or
        localhost (\texttt{N2P2\_COMPUTER}, code labels via environment variables);
  \item \texttt{examples/1.Boron/aiida-n2p2\_demo.py} --- boron example with
        \texttt{Int(5)} and \texttt{retrieve\_trajectory = Bool(False)};
  \item \texttt{examples/3.Restart/restart\_train\_demo.py} --- manual restart:
        two WorkChain submits (default 40~+~40 epochs; Al inputs from
        \texttt{examples/2.HPC/});
  \item \texttt{examples/3.Restart/auto\_restart\_dahu.py} --- automatic restart
        on Dahu/OAR (\texttt{pr-diamond}, configurable walltime per segment);
  \item \texttt{examples/3.Restart/auto\_restart\_train\_demo.py} --- automatic
        restart smoke test (any scheduler; localhost walltime may not trigger
        multiple segments).
\end{itemize}

Restart examples require the full Al \texttt{input.data} from
\texttt{examples/2.HPC/}; the minimal files under \texttt{examples/3.Restart/} are
for unit tests only.
